\subsection{Profile Distribution}\label{sec:res-profile}
Learning a new fabric from $k=10$ frames, including both LOO calibrations, took 0.55--0.96\,s for the station model and 0.17\,s for PatchCore-lite (3.6\,s on the very first fabric, which includes loading the ResNet-18 weights). Publishing the 6.1\,kB station profile as a retained message delivered it to every subscribed station in 10--15\,ms over the Wi-Fi link and in 93\,ms over the constrained link, independently of whether 1, 8, or 32 stations subscribed, because the broker fans the message out once per station link. Distributing a PatchCore-lite profile instead (1.25\,MB) would take 0.52\,s and 10.1\,s, respectively. A fabric change therefore costs the plant a few seconds of operator time for capturing frames, not a model-deployment step.

\subsection{Robustness to Imaging Conditions}\label{sec:res-robust}
Table~\ref{tab:robust} applies global brightness changes, defocus blur, and sensor noise to the test frames at the $256\times256$ working resolution, with models fitted on clean shots. All detectors are insensitive to $\pm$20\,\% brightness in terms of AUROC (at most $-0.009$): a global gain shifts the scores of good and defective frames alike, so their ranking is preserved, although the absolute LOO threshold may then be off. Blur ($\sigma=2$\,px) costs 0.08--0.12 AUROC for all three detectors because it removes the fine texture that both Gabor energies and early ResNet layers rely on. Noise ($\sigma=10$ grey levels) affects the station model most ($-0.128$), since its finest Gabor scale ($\lambda=4$\,px) responds to pixel noise, and PatchCore-lite least ($-0.048$). In practice, a change of focus, lighting or camera should be followed by re-learning the fabric, which, as shown above, takes seconds and needs no labels; the logged score distribution of recent OK frames can be used to trigger it automatically.

\begin{table*}[t]
\centering
\caption{Robustness of image AUROC to imaging changes applied to the test frames only ($k=20$, mean over five categories; change w.r.t.\ clean in parentheses).}
\label{tab:robust}
\footnotesize
\begin{tabular}{@{}lccccc@{}}
\toprule
Model & none & bright-20\% & bright+20\% & blur($\sigma$=2) & noise($\sigma$=10) \\
\midrule
GG & 0.872 & 0.879 (+0.007) & 0.864 (-0.008) & 0.772 (-0.099) & 0.744 (-0.128) \\
PC & 0.957 & 0.955 (-0.002) & 0.957 (+0.000) & 0.878 (-0.079) & 0.909 (-0.048) \\
PD & 0.904 & 0.895 (-0.009) & 0.903 (-0.001) & 0.783 (-0.121) & 0.850 (-0.054) \\
\bottomrule
\end{tabular}
\end{table*}

